\thispagestyle{empty}\vspace*{\fill}
\begin{minipage}{128mm}\fontsize{9}{13}\selectfont\color{otMuted}
{\otSemi\color{otInk}Mathématiques · Troisième · Brevet 2027}\par\medskip
Conception et mise en page : collection OptiTAB.\par
270 exercices de chapitre, 18 défis et deux brevets blancs originaux, corrigés.\par
Charte graphique de la collection OptiTAB de Terminale, Première et Seconde.\par\bigskip
Édition du 21 octobre 2026 · France · Année scolaire 2026-2027.\par\bigskip
Le livre suit les attendus de Troisième actuellement applicables et réactive
les connaissances du cycle 4 nécessaires. Le nouveau programme publié en 2026
ne s’appliquera en Troisième qu’en 2028-2029. Le format des brevets blancs
s’appuie sur la définition officielle de septembre 2026, applicable à compter
de la session 2027.\par\bigskip
L’ordre des chapitres, les exercices et les barèmes d’entraînement sont des
choix éditoriaux. Les défis sont facultatifs. Les situations et données
numériques sont construites pour l’apprentissage ; elles ne constituent pas
des données économiques ou scientifiques relevées.\par\bigskip
Les schémas de principe ne permettent pas de mesurer une longueur inconnue.
Les coordonnées et graduations explicites font foi pour les graphiques.
Sauf demande d’arrondi, conserver les résultats exacts.\par\bigskip
© OptiTAB, 2026.\par\otSite
\end{minipage}\clearpage
